% Registered A.53: independent root and DE full proof/source/code audits PASS.
\subsection{A fixed-subgroup enlargement on every cover compactification}

\begin{statement}{A.53}{Proved}{Every cover compactification admits a torsion enlargement with full fixed inertia}
Let \(Y\to U\) be any connected finite etale cover from Task A.5,
with proper compactification \(C\), degree \(d\), proper genus
\(g\), and \(r\) cusps. Let \(q:C\to E\) be its extended
proper map. There are infinitely many rational primes
\(p\equiv3\pmod4\) such that
\[
       Z_p=C\setminus q^{-1}(E[p]),\qquad
                    Z_p\xrightarrow{[p]q}U             \tag{A.53.1}
\]
is a connected finite etale cover of degree \(p^2d\), proper genus
\(g\), and \(p^2d-2g+2\) cusps, and its full geometric
pro-\(\ell\) fixed group is the selected cusp inertia for every
coefficient prime \(\ell\), every cusp packet, and every indicated
open derived subgroup.

If \(Y\) has cuspidal procyclic normalizer recognition at a given
coefficient prime, then every further puncturing of \(Y\), in
particular \(Z_p\), has that recognition at the same prime.
Thus any previously proved all-prime normalizer theorem for \(Y\)
also holds for (A.53.1).

The comparison \(Z_p\hookrightarrow Y\) is an open immersion;
the new map to \(U\) is the composite with \([p]\). This construction
asserts neither finite refinement over the original map to \(U\)
nor lifting of fixed elements from \(Y\). In particular it does
not prove the universal assertion of Task A.5.
\end{statement}

\begin{proof}
Enlarge the number field so that all original cusps are rational.
Their proper genus satisfies \(g\ge1\), since \(q\) is a finite
map to an elliptic curve. The map is ramified only above the origin,
and its full origin fibre is precisely the cusp set. Let \(e_c\)
be its ramification index at cusp \(c\). Thus
\[
                   \sum_c e_c=d,\qquad
                         2g-2=\sum_c(e_c-1)=d-r.    \tag{A.53.2}
\]

\paragraph{Choose an actual good tame model.}
Spread out the proper curve, the finite morphism, and the finitely
many cusp sections over an open subset of the integer ring. After
removing finitely many primes, both proper curves are smooth, the
cusp sections are pairwise disjoint, \(q\) is finite flat of degree
\(d\), and it is etale off the origin section. Moreover the equality
of Cartier divisors
\(q^*(O)=\sum_c e_c(c)\) holds on these models.
These assertions can be obtained by clearing denominators in finite
presentations for the morphism and the displayed divisor equality,
and then removing the closed nonsmooth, nonflat, and unwanted
ramification loci, which have empty generic fibre. In particular
only finitely many rational primes have been excluded. The
finite-bad-reduction assertion is also
\cite[Proposition 10.1.21(a), pp.~462--463]{fghLiu}.

Choose a rational prime \(p\equiv3\pmod4\) outside that finite set
and with \(p>d\). There are infinitely many choices. This inequality
makes every index \(e_c\) prime to \(p\); it avoids inferring
individual tame indices merely from a degree coprimality. At a place
above \(p\), \(E:y^2=x^3-x\) has supersingular good reduction;
see \cite[Theorem V.4.1(a) and Example V.4.5, pp.~148 and 152]{fghSilvermanAEC}.
After a further finite extension make all points of \(E[p]\) and
their inverse images rational. Good reduction and the stated finite
map persist. Every point of \(E[p]\) specializes to the origin,
since a supersingular elliptic curve has no nonzero geometric
\(p\)-torsion in its special fibre.

\paragraph{Every cusp has a companion in its own rational cluster.}
Fix an original cusp \(c\). In the completed strict henselian
neighborhoods of the source and target sections choose parameters
\(u,t\). The Cartier-divisor equality gives
\[
                              t=u^{e_c}h(u),
                                                               \tag{A.53.3}
\]
where \(h\) is a unit. Since \(e_c\) is invertible in the residue
field, extract an \(e_c\)-th root of this unit and change \(u\).
The map is then exactly \(t=u^{e_c}\). Each of the \(p^2-1\)
nonzero characteristic-zero points of \(E[p]\) has nonzero parameter
of positive valuation. Its \(e_c\) distinct roots all reduce to
\(c\). The origin itself has one source point in this neighborhood.
Consequently the number of marked points specializing to \(c\) is
\[
                              1+e_c(p^2-1)\ge p^2.  \tag{A.53.4}
\]
There are no other reduction classes: all target points specialize
to the origin, and the special origin fibre is the disjoint set of
original cusp points with its indicated multiplicities. These local
counts sum to all points of \(q^{-1}(E[p])\).

The proper source model is smooth. Separating each of these clusters
by blowups, after finite extension if necessary, creates rational
marked trees attached at the corresponding original smooth source
point. For any chosen new cusp, another mark lies in its cluster by
(A.53.4), and their entire joining path, including both marked
components, is rational. The proper model has compact type because
it is smooth; inserting these rational trees preserves that property.
This is exactly the pointed geometric hypothesis of A.41--A.42.
Those statements give the full integral fixed-inertia assertion at
every coefficient prime, including \(\ell=p\), for every packet
and every indicated open derived group. No restriction on the
\(p\)-rank or endomorphism algebra of \(J_C\) is required here.

The characteristic-zero map \([p]q\) has degree \(p^2d\) and is
unramified after removal of its full origin fibre. The new affine
curve is geometrically connected because it is an open subset of
the same smooth connected proper curve. Its cusp count is
\[
 \sum_c\bigl(1+e_c(p^2-1)\bigr)
       =r+(p^2-1)d=p^2d-2g+2.                       \tag{A.53.5}
\]
This also verifies its Euler characteristic \(-p^2d\).

\paragraph{Normalizer recognition survives further puncturing.}
More generally let \(Z\hookrightarrow Y\) be an open immersion of
smooth affine curves with the same proper compactification. Suppose
that \(Y\) satisfies the pro-\(\ell\) normalizer criterion. If a
noncuspidal normalized procyclic subgroup existed in the geometric
pro-\(\ell\) group of \(Z\), A.11 would give it nonzero image in
the proper Tate homology of that common compactification. Its image
under the actual morphism \(Z\to Y\) is therefore nontrivial and
has the same nonzero proper-homology image. The arithmetic normalizer
maps to an arithmetic normalizer with open projection. The criterion
on \(Y\) puts this image in cusp inertia, whose proper-homology
image is zero. This contradiction proves the assertion.
The converse direction, that subgroups of cusp inertia have open
arithmetic normalizer after finite field restriction, is the usual
cusp decomposition-group inclusion.

Finally, the actual comparison to \(Y\) is the open immersion, which
induces a surjection on geometric pro-prime fundamental groups.
The fixed-group functor is not right exact, so the fixed theorem
for \(Z_p\) is not being pushed down to \(Y\). The map (A.53.1)
is a new covering map to \(U\), rather than a factorization through
\(Y\to U\) as a finite cover of that same base.
\end{proof}
